\documentclass[aspectratio=169]{beamer}

\usetheme{Madrid}
\usecolortheme{default}

\usepackage[T1]{fontenc}
\usepackage[utf8]{inputenc}
\usepackage{lmodern}
\usepackage{graphicx}
\usepackage{hyperref}

\title{Artemis II Mission Timeline Analysis}
\author[Group XXX]{
  Group number: XXX \\
  Student 1 \\
  Student 2 \\
  Student 3
}
\institute{Modelling Complex Systems, Spring 2026}
\date{Lab 0 Practice Presentation}

\begin{document}

\begin{frame}
  \titlepage
\end{frame}

\begin{frame}{Question}
  State the question you want to answer with the dataset.

  \vspace{0.5em}
  Example:
  \begin{block}{Example question}
    How do the spacecraft's distance to Earth and distance to the Moon change during the main stages of Artemis II?
  \end{block}
\end{frame}

\begin{frame}{Data and code}
  Briefly explain:
  \begin{itemize}
    \item which file you used,
    \item which variables you plotted,
    \item which Python script produced the figures,
    \item how missing values were handled.
  \end{itemize}
\end{frame}

\begin{frame}{Main figure}
  \begin{figure}
    \centering
    % \includegraphics[width=0.8\textwidth]{distancePlot.pdf}
    \caption{Replace with your main plot and caption.}
  \end{figure}
\end{frame}

\begin{frame}{Interpretation}
  Summarize the main message of the figure.

  \vspace{0.5em}
  Possible points:
  \begin{itemize}
    \item outbound travel,
    \item lunar flyby,
    \item return trajectory,
    \item key event markers.
  \end{itemize}
\end{frame}

\begin{frame}{Takeaway}
  End with two or three short points:
  \begin{itemize}
    \item the main result,
    \item one data or plotting choice that mattered,
    \item one thing your group learned from the workflow.
  \end{itemize}
\end{frame}

\end{document}
